\documentclass[12pt,a4paper]{article}
\usepackage[utf8]{inputenc}
\usepackage[turkish]{babel}
\usepackage{amsmath}
\usepackage{graphicx}
\usepackage{hyperref}
\usepackage{listings}
\usepackage{xcolor}
\usepackage{enumitem}
\usepackage{geometry}
\geometry{a4paper, margin=1in}

% Define colors for code listings
\definecolor{codegreen}{rgb}{0,0.6,0}
\definecolor{codegray}{rgb}{0.5,0.5,0.5}
\definecolor{codepurple}{rgb}{0.58,0,0.82}
\definecolor{backcolour}{rgb}{0.95,0.95,0.92}

\lstdefinestyle{mystyle}{
    backgroundcolor=\color{backcolour},   
    commentstyle=\color{codegreen},
    keywordstyle=\color{magenta},
    numberstyle=\tiny\color{codegray},
    stringstyle=\color{codepurple},
    basicstyle=\ttfamily\footnotesize,
    breakatwhitespace=false,         
    breaklines=true,                 
    captionpos=b,                    
    keepspaces=true,                 
    numbers=left,                    
    numbersep=5pt,                  
    showspaces=false,                
    showstringspaces=false,
    showtabs=false,                  
    tabsize=2
}
\lstset{style=mystyle}

\begin{document}
\shorthandoff{=}

\begin{titlepage}
    \centering
    \vspace*{1cm}
    
    {\Large\textbf{ANKARA UNIVERSITESI}}\\[0.3cm]
    {\large Muhendislik Fakultesi}\\[0.3cm]
    {\large Bilgisayar Muhendisligi Bolumu}\\[1cm]

    \begin{figure}[h]
    \centering
    \vspace{3pt}
    \includegraphics[width=5.0cm]{logo.png}
    \end{figure}

    \begin{tabular}{ll}
        \textbf{Ogrenci:} & Ahmet Akgun \\
        \textbf{Ogrenci No:} & 22290602 \\
        \textbf{Danisman:} & Enver Bagci \\
    \end{tabular}\\[1.5cm]
    
    {\Large \textbf{BLM4522 Ag Tabanli Paralel Dagitim Sistemleri}}\\[0.5cm]
    {\large \textbf{Proje Adi:} Veritabani Yedekleme ve Otomasyon Calismasi}
    
    \vfill
    
    \begin{itemize}[label=\textbullet]
        \item \textbf{GitHub Depo Linki:} \url{https://github.com/aakgunnn}
        \item \textbf{Haftalık Video Takip Bağlantısı:} \url{https://www.youtube.com/@ahmetakgun359}
    \end{itemize}

\end{titlepage}

\tableofcontents
\newpage

\section{Giriş}
Bu proje kapsamında, kurumsal düzeyde bir veritabanının günlük ve operasyonel yedekleme (backup) süreçlerinin nasıl otomatikleştirildiği uygulamalı olarak incelenmiştir. Geleneksel SQL Server ortamlarındaki "SQL Server Agent" mekanizmasının işlevleri, PostgreSQL veritabanı yönetim sisteminde ve Windows İşletim Sistemi mimarisinde harmanlanarak esnek bir altyapı oluşturulmuştur.

\section{Aşama 1: Örnek Veritabanı (Northwind)}
Projeyi en verimli ve gerçekçi şartlarda test edebilmek amacıyla, GitHub platformu üzerinden dünyaca ünlü, tabloları tam dolu olan \textbf{Northwind} \texttt{(northwind-pg)} açık kaynak örnek veritabanı diske indirilmiştir. İndirildiği gibi hedef PostgreSQL sunucusuna dahil edilerek \texttt{yedek\_demo\_db} adlı veritabanı üzerinden test senaryolarına geçilmiştir.

\section{Aşama 2: Yedekleme Otomasyonu ve Script (PowerShell)}
PostgreSQL tabanlı sistemlerde yedek işlemi temel olarak \texttt{pg\_dump} komut aracına dayanır. Bu projede rutin bir sistem yöneticisinin elle yazacağı komutları dışlayan, hataları yöneten ve raporlayan bir \texttt{backup\_manager.ps1} isimli PowerShell betiği (script) kodlanmıştır.

\begin{lstlisting}[language=bash, caption=pg\_dump Otomasyon Mantığı]
pg_dump.exe -U postgres -d yedek_demo_db -F c -f "Backups\backup_2026...backup"
\end{lstlisting}

Yukarıda görüldüğü gibi sisteme \textbf{-F c (Custom Format)} parametresi geçilerek veritabanı yedeğinin sıkıştırılmış olan spesifik \texttt{.backup} formunda tutulması zorunlu kılınmıştır. Bu sayede sunucu diskinde oluşacak alan tüketimi minimalize edildiği gibi kurtarma (Restore) işlemleri çok daha kontrollü yapılabilmektedir.

\section{Aşama 3: Denetim ve Raporlama (Audit Logs)}
Kurumsal veri yönetimi, alınan güvenlik yedeklerinin doğrulanmasını katı standartlara bağlar. Yazılan otomasyon içerisinde yer alan analiz mekanizmaları kullanılarak yedeklerin başarım durumlarına dair özet veriler toplanmıştır:

\begin{itemize}
    \item İşlemin tam olarak kaç saniye (veya milisaniye) sürdüğü (\texttt{Measure-Command} sınıfı kullanımı).
    \item Hedef .backup dosyasının MB olarak net ağırlığı.
    \item Zaman damgası (Timestamp) bazlı sonuç değerlendirmesi.
\end{itemize}

Otomatik betik bu metrikleri harmanlayıp anlık olarak \texttt{Logs/backup\_audit\_log.txt} isminde özel bir izleme dosyasına yazmaktadır:
\begin{lstlisting}[caption=Audit Log Çıktısı Örneği]
[2026-04-07 12:47:05] [SUCCESS] DB: yedek_demo_db | ZAMAN: 1.11 sn | BOYUT: 0.06 MB | DOSYA: backup_2026.backup
[2026-04-07 12:47:21] [ERROR] DB: yedek_demo_db | DETAY: baglanti hatasi
\end{lstlisting}

\section{Aşama 4: Otomatik Uyarı Mekanizması}
Herhangi bir port kesintisi, sunucu donanım arızası ya da yetki kaybı dolayısıyla yedekleme işlemi \textbf{başarısız} olduğunda (Exception atıldığında), acil bir reaksiyon gösterilmesi kritik önem taşır.

Scripte eklenen yapı sayesinde, otomasyon betiği anında bir sanal Mail tetiklemektedir. Gerçek bir SMTP relay'in yorulmasına ihtiyaç duyulmaksızın proje sürecinde uyarıların ulaştığını görmek adına simüle (Mock) bir işlev kodlanmıştır:
\begin{lstlisting}[caption=Mail Simülasyon Çıktısı]
[EMAIL SIMULATION] TO: dbadmin@company.com 
SUBJECT: [CRITICAL] Yedekleme Basarisiz! 
BODY: Veritabani yedeklemesi tamamlanamadi. Lutfen loglari inceleyin.
\end{lstlisting}
Yazılan bu mantık sayesinde mail gönderim portları açıldığı gün operasyonel kod saniyeler içerisinde entegre edilebilecektir.

\section{Aşama 5: Görev Zamanlayıcısı (Task Scheduler) Entegrasyonu}
Tüm bu kusursuz betiğin insan eli değmeden otomatik koşabilmesi için SQL Server Agent yeteneğinin Windows karşılığı olan Gelişmiş Görev Zamanlayıcısı (\texttt{schtasks}) devrede alınmıştır.

\begin{lstlisting}[language=bash, caption=Zamanlanmış PowerShell Görevinin Entegre Edilmesi]
schtasks /create /tn "PostgreSQL_Nightly_Backup" /tr "powershell.exe -ExecutionPolicy Bypass -WindowStyle Hidden -File backup_manager.ps1" /sc daily /st 02:00
\end{lstlisting}

Geceleri sistem saat \textbf{02:00}'ı vurduğunda yukarıdaki entegrasyon aracılığı ile yedek mekanizması sıfır işlemci penceresi ile şeffaf çalışacak (Hidden), yedeğini kaydedecek ve logu sisteme bırakacaktır.

\section{Sonuç}
Bu projede PostgreSQL veritabanları üzerinde yedekleme otomasyonunun temelleri modern Windows işletim sistemleriyle haberleştirilerek incelenmiştir. Yedekleme hızı, disk optimizasyonu, otomatik denetim, raporlama dökümü ve uyarı (notification) altyapısı bir araya getirilerek eksiksiz bir mimari entegrasyonu simüle edilmiştir. 

\end{document}
